\documentclass[11pt]{article}

% --- Packages ---
\usepackage[utf8]{inputenc}
\usepackage[T1]{fontenc}
% \usepackage{lmodern}
% \usepackage{microtype}
\usepackage{geometry}
\geometry{margin=1in}

\usepackage{amsmath,amssymb,amsthm,mathtools}
\usepackage{stmaryrd}
\usepackage{enumitem}
\usepackage{tikz-cd}
\usepackage{hyperref}
\hypersetup{colorlinks=true,linkcolor=blue,citecolor=blue,urlcolor=blue}
\usepackage{mdframed}
\usepackage{booktabs}

% --- Theorem envs ---
\theoremstyle{definition}
\newtheorem{definition}{Definition}[section]
\newtheorem{remark}[definition]{Remark}
\newtheorem{example}[definition]{Example}
\theoremstyle{plain}
\newtheorem{proposition}[definition]{Proposition}
\newtheorem{theorem}[definition]{Theorem}
\newtheorem{lemma}[definition]{Lemma}
\newtheorem{corollary}[definition]{Corollary}

% --- Macros ---
\newcommand{\Pow}{\mathcal P}
\newcommand{\Ext}{\mathrm{Ext}}
\newcommand{\Int}{\mathrm{Int}}
\newcommand{\Intt}{\mathrm{Int}}
\newcommand{\Hit}{\mathrm{Hit}}
\newcommand{\Sel}{\mathrm{Sel}}
\newcommand{\cl}{\mathrm{cl}}
\newcommand{\Opp}{\mathsf{op}}
\newcommand{\forces}{\Vdash}
\newcommand{\ltimesrel}{\ltimes}
\newcommand{\covers}{\triangleleft}
\newcommand{\Sub}{\mathrm{Sub}}
\newcommand{\Sh}{\mathrm{Sh}}
\newcommand{\Split}{\mathsf{Split}}

\title{Positive Topology and Feasible Refinement: \\ Forcing Matrices, Positivity, and Information}
\author{Mirco A. Mannucci, Giovanni Sambin}
\date{\today}

\begin{document}
\maketitle

\begin{center}
\textbf{Keywords:} positive topology, formal topology, pointfree topology, locales and frames, overt locales (overtness), positivity relation $a\ltimes U$, Galois connections and adjunctions, forcing semantics (forcing matrices), closed sets without complements, information refinement, game semantics, constructive mathematics, resource-/cost-aware semantics

\medskip
\textbf{MSC 2020:} 06D22 (Frames, locales), 03F65 (Other constructive mathematics), 06D20 (Heyting algebras), 18B25 (Toposes), 18A40 (Adjunctions)

\medskip
\textbf{ACM CCS:} Theory of computation $\to$ Logic; Semantics and reasoning; Program semantics; Type theory
\end{center}


\begin{abstract}
Starting from a point--basic-open forcing relation $x\forces a$ between a set of points/models $X$
and basic opens/generators $S$, we obtain the canonical Galois adjunction
$\Ext \dashv \Intt$ between $\Pow(S)$ and $\Pow(X)$ encoding the ``open/locale'' (universal, cover/refinement) content~\cite{MacLaneMoerdijk1992}.
Less standard, but equally canonical, is a second adjunction $\Hit \dashv \Sel$ encoding the ``closed/positive''
(existential, probe/meet) content; this yields a direct origin for Sambin-style positivity $a\ltimesrel U$
and ``closed without complements'' as fixed points of the contractive operator $J=\Hit\circ\Sel$~\cite{Sambin-OUP-PositiveTopology,CoquandSambinSmithValentini2003}.
We isolate a reconstruction principle: the forcing matrix is recovered uniquely from either residual pair via singleton tests.
Costs enter by replacing Boolean truth values with values in an ordered cost structure, relativizing point vs.\ point-free
(weighted points, thresholds, and enriched closures).

Treating positivity as \emph{primitive} exposes two operational readings~\cite{Valentini2012,CirauloVickers2016}:
(i) an information-theoretic view, where covers are feasible refinements of partial information and positivity is witnessed feasibility, and
(ii) a game view, where positivity corresponds to a winning strategy in a cover-refinement game.
Together they make the ``two interlaced adjunctions'' explicit and prepare the ground for graded/resource-bounded truth.

We illustrate these abstract notions through concrete examples from medical diagnosis, legal reasoning, and AI systems,
showing how the framework provides a foundation for grounded, explainable, and resource-aware inference.
\end{abstract}

\tableofcontents

% =========================================================
\section{Introduction: Why Positive Topology?}

Before diving into the formal machinery, we motivate the key ideas through a concrete scenario.

\subsection{A motivating example: Medical diagnosis}

Consider the task of diagnosing a patient based on observed symptoms.

\begin{itemize}
  \item $X$ = patients (the ``points'' or ``states of the world'')
  \item $S$ = symptoms, tests, conditions (the ``observables'' or ``basic opens'')
  \item $x \forces a$ means ``patient $x$ exhibits symptom/condition $a$''
\end{itemize}

A doctor observes that a patient has fever, cough, and fatigue. What can we infer?

\paragraph{The cover relation: semantic entailment.}
Certain conditions \emph{imply} certain symptoms. For instance:
\[
\mathsf{pneumonia} \covers \{\mathsf{fever}, \mathsf{cough}, \mathsf{chest\_pain}, \mathsf{infiltrates}\}
\]
This means: every patient with pneumonia exhibits at least one of fever, cough, chest pain, or lung infiltrates.
More generally, $a \covers U$ says: \emph{whenever $a$ holds, something in $U$ holds}.

This is the \textbf{universal/refinement} side: covers encode what implies what.

\paragraph{The positivity relation: witnessed consistency.}
Now consider the question: is the symptom cluster $\{\mathsf{fever}, \mathsf{cough}, \mathsf{fatigue}\}$ a \emph{realizable} profile?
That is, does there exist a patient whose \emph{complete} symptom profile is exactly this set?

We write:
\[
\mathsf{fever} \ltimesrel \{\mathsf{fever}, \mathsf{cough}, \mathsf{fatigue}\}
\]
to mean: there exists a patient with fever whose entire observable profile is contained in $\{\mathsf{fever}, \mathsf{cough}, \mathsf{fatigue}\}$.

This patient is a \textbf{witness} --- the symptom cluster is realizable, not a contradiction.

Contrast with:
\[
\mathsf{fever} \ltimesrel \{\mathsf{hypothermia}\}
\]
This is \emph{false}: no patient has both fever and only hypothermia. The combination is inconsistent.

This is the \textbf{existential/positivity} side: positivity encodes what is consistent, what has witnesses.

\paragraph{Why both matter.}
\begin{itemize}
  \item \textbf{Covers alone} tell us what implies what, but not what's realizable.
  \item \textbf{Positivity alone} tells us what's consistent, but not what follows from what.
  \item \textbf{Together} they give us: grounded inference (every conclusion has a witness) with semantic structure (implications are tracked).
\end{itemize}

\paragraph{The compatibility axiom.}
The key interaction between covers and positivity is:
\[
(a \covers U) \land (a \ltimesrel V) \implies \exists u \in U.\; u \ltimesrel V
\]
In medical terms: if pneumonia implies \{fever, cough, infiltrates\}, and there exists a pneumonia patient confined to profile $V$,
then there exists a fever patient (or cough patient, or infiltrates patient) also confined to $V$.

\textbf{Refinement preserves witnesses.} This is the soundness condition that makes the framework coherent.

\subsection{The general pattern}

The medical example instantiates a universal pattern:

\begin{center}
\begin{tabular}{lll}
\toprule
\textbf{Domain} & \textbf{Points $X$} & \textbf{Observables $S$} \\
\midrule
Medicine & Patients & Symptoms, tests, conditions \\
Law & Cases & Evidence, claims, precedents \\
Programming & Program states & Properties, assertions, tests \\
Finance & Portfolios & Risk factors, exposures \\
Science & Hypotheses & Experiments, observations \\
AI agents & Traces/executions & Constraints, evaluations \\
\bottomrule
\end{tabular}
\end{center}

This operational reading follows the constructive-positivity tradition initiated by Sambin~\cite{Sambin-OUP-PositiveTopology}
and formalized via inductively generated formal topologies~\cite{CoquandSambinSmithValentini2003}.
It abstracts the familiar locale/sheaf perspective of Mac~Lane--Moerdijk~\cite{MacLaneMoerdijk1992} down to the level of the forcing matrix we use in code.

In each domain:
\begin{itemize}
  \item $a \covers U$ encodes \emph{semantic entailment}: $a$ implies something in $U$
  \item $a \ltimesrel U$ encodes \emph{witnessed consistency}: there's a state with $a$ confined to $U$
  \item The forcing $x \forces a$ connects the abstract observables to concrete states
\end{itemize}

The rest of this paper develops these ideas formally, showing how both structures arise canonically from a single forcing relation,
and how adding costs/resources yields a budget-aware semantics suitable for AI engineering.

% =========================================================
\section{The starting datum: a forcing/incidence relation}

\begin{definition}[Forcing table]
Fix sets $X$ (points/models) and $S$ (basic opens/generators).
A \emph{forcing relation} is a subset $R\subseteq X\times S$ written
\[
x\forces a \quad\text{for }(x,a)\in R.
\]
\end{definition}

\begin{remark}[Interpretation]
Think of $a\in S$ as a basic open (or observable, or test) and $x\in X$ as a point (or state, or model).
Then $x\forces a$ reads: ``$x$ lies in the basic open $a$'' or ``$x$ satisfies observable $a$''.
No topology is assumed at the outset: we start from the incidence matrix.
\end{remark}

\begin{example}[Medical diagnosis]\label{ex:medical-forcing}
Let $X$ be a set of patients and $S$ a set of symptoms/conditions.
Then $x \forces a$ means ``patient $x$ exhibits symptom $a$''.
The forcing relation is the patient-symptom incidence matrix.
\end{example}

\begin{example}[Program verification]
Let $X$ be program states and $S$ be assertions/properties.
Then $x \forces a$ means ``state $x$ satisfies property $a$''.
\end{example}

\begin{remark}[Boolean matrix viewpoint]
The relation $R\subseteq X\times S$ is a Boolean-valued matrix
$R:X\times S\to \{0,1\}$.
Much of what follows is ``linear algebra over $\{0,1\}$'': joins are unions, meets are intersections.
This viewpoint becomes crucial when we replace $\{0,1\}$ by a cost/weight structure.
\end{remark}

% =========================================================
\section{The first adjunction: opens via $\Ext \dashv \Intt$}\label{sec:adj1}

\subsection{Direct image and the universal residual}

\begin{definition}[Extension and interior-kernel]
Define monotone maps
\[
\Ext:\Pow(S)\to\Pow(X),\qquad \Ext(U)=\{x\in X:\exists a\in U,\ x\forces a\},
\]
\[
\Intt:\Pow(X)\to\Pow(S),\qquad \Intt(A)=\{a\in S:\forall x\in X,\ (x\forces a\Rightarrow x\in A)\}.
\]
\end{definition}

\begin{example}[Medical: Extension and Interior]
Continuing Example~\ref{ex:medical-forcing}:
\begin{itemize}
  \item $\Ext(\{\mathsf{fever}, \mathsf{cough}\})$ = patients who have fever OR cough (or both)
  \item $\Intt(A)$ where $A$ = ``patients over 65'' gives: symptoms that \emph{only} occur in patients over 65
\end{itemize}
\end{example}

\begin{proposition}[Galois adjunction I]
For all $U\subseteq S$ and $A\subseteq X$,
\[
\Ext(U)\subseteq A \quad\Longleftrightarrow\quad U\subseteq \Intt(A).
\]
Equivalently, $\Ext\dashv \Intt$.
\end{proposition}

\begin{proof}
($\Rightarrow$) If $a\in U$, then $\Ext(\{a\})\subseteq \Ext(U)\subseteq A$, hence $a\in \Intt(A)$.
($\Leftarrow$) If $x\in \Ext(U)$, choose $a\in U$ with $x\forces a$. Since $a\in\Intt(A)$, we have $x\in A$.
\end{proof}

\begin{remark}[Universal nature]
$\Intt(A)$ is defined by a $\forall$-condition; this is the ``universal/refinement'' side.
This is exactly why the cover relation induced by forcing is a universal statement.
\end{remark}

\subsection{Cover and the point-free (locale/formal topology) side}

\begin{definition}[Induced preorder and cover]
Define a preorder on $S$ by
\[
a\le b \quad:\Longleftrightarrow\quad \forall x,\ (x\forces a\Rightarrow x\forces b),
\]
and define a cover relation $\covers$ by
\[
a\covers U \quad:\Longleftrightarrow\quad \forall x,\ (x\forces a \Rightarrow \exists b\in U,\ x\forces b).
\]
\end{definition}

\begin{example}[Medical: Cover as symptom entailment]\label{ex:medical-cover}
In the medical setting, $a \covers U$ means: every patient with condition $a$ exhibits at least one symptom in $U$.
For instance:
\begin{align*}
\mathsf{flu} &\covers \{\mathsf{fever}, \mathsf{cough}, \mathsf{fatigue}, \mathsf{myalgia}\} \\
\mathsf{bacterial\_infection} &\covers \{\mathsf{fever}, \mathsf{elevated\_WBC}\}
\end{align*}
These are \emph{diagnostic implications}: the disease implies certain observable symptoms.
\end{example}

\begin{remark}[Extensional reading]
Equivalently:
\[
a\le b \iff \Ext(\{a\})\subseteq \Ext(\{b\}),\qquad
a\covers U \iff \Ext(\{a\})\subseteq \Ext(U).
\]
So $\covers$ is literally ``inclusion in a union''.
\end{remark}

\subsection{Extensional topology on points}

\begin{definition}[Generated opens on $X$]
A subset $A\subseteq X$ is \emph{extensional open} (generated by $S$) if
\[
A=\Ext(\Intt(A)).
\]
\end{definition}

\begin{remark}
The operator $c:=\Ext\circ\Intt$ is an \emph{interior} operator on $\Pow(X)$: it extracts the largest extensional open contained in a given subset.
(Dually, $j:=\Intt\circ\Ext$ is a closure/nucleus-like operator on $\Pow(S)$.)
Extensional opens are exactly the fixed points of $c$: ``those subsets describable by generators''.
\end{remark}

\begin{definition}[Extensional equivalence of points]
Define $x\sim y$ if the induced neighborhood sets agree:
\[
x\sim y \quad:\Longleftrightarrow\quad \forall a\in S,\ (x\forces a \Leftrightarrow y\forces a).
\]
\end{definition}

\begin{example}[Medical: Indistinguishable patients]
Two patients $x \sim y$ are \emph{symptomatically indistinguishable}: they exhibit exactly the same symptoms.
From a diagnostic standpoint, they are equivalent --- the same inferences apply to both.
\end{example}

\begin{remark}
Point-free semantics does not distinguish points inside the same $\sim$-class.
This is the basic ``quotient-to-locale'' phenomenon: locale data sees neighborhood behavior, not point multiplicity.
\end{remark}

% =========================================================
\section{The second adjunction: positivity via $\Hit \dashv \Sel$}\label{sec:adj2}

The locale/cover side is universal; ``closed without complements'' needs an existential probe:
``does a basic open \emph{hit} a closed thing?''.
This is the conceptual origin of Sambin-style positivity~\cite{Sambin-OUP-PositiveTopology,Valentini2012,CirauloVickers2016}.

\subsection{Neighborhood sets and constraint-selected points}

\begin{definition}[Neighborhood set]
For $x\in X$, define its neighborhood set (or \emph{observable profile}) by
\[
N(x):=\{a\in S:\ x\forces a\}\subseteq S.
\]
\end{definition}

\begin{example}[Medical: Symptom profile]
For a patient $x$, the neighborhood set $N(x)$ is their \emph{complete symptom profile}:
the set of all symptoms/conditions they exhibit.
\end{example}

\begin{definition}[Selection by a constraint on neighborhoods]
Define
\[
\Sel:\Pow(S)\to\Pow(X),\qquad
\Sel(U):=\{x\in X:\ N(x)\subseteq U\}.
\]
Thus $\Sel(U)$ consists of points all of whose basic neighborhoods lie in $U$.
\end{definition}

\begin{example}[Medical: Patients confined to a symptom cluster]
$\Sel(\{\mathsf{fever}, \mathsf{cough}, \mathsf{fatigue}\})$ = patients whose \emph{entire} symptom profile
is contained in $\{\mathsf{fever}, \mathsf{cough}, \mathsf{fatigue}\}$.

These are patients who have \emph{only} fever, cough, and/or fatigue --- no other symptoms.
\end{example}

\subsection{The existential ``hit'' map}

\begin{definition}[Hit]
Define
\[
\Hit:\Pow(X)\to\Pow(S),\qquad
\Hit(C):=\{a\in S:\exists x\in C,\ x\forces a\}.
\]
\end{definition}

\begin{example}[Medical: Symptoms witnessed in a cohort]
If $C$ = patients in a clinical trial, then $\Hit(C)$ = symptoms exhibited by at least one patient in the trial.
\end{example}

\begin{proposition}[Galois adjunction II]
For all $C\subseteq X$ and $U\subseteq S$,
\[
\Hit(C)\subseteq U \quad\Longleftrightarrow\quad C\subseteq \Sel(U).
\]
Equivalently, $\Hit\dashv \Sel$.
\end{proposition}

\begin{proof}
($\Rightarrow$) Let $x\in C$. If $a\in N(x)$ then $a\in\Hit(C)\subseteq U$. Hence $N(x)\subseteq U$ and $x\in\Sel(U)$.

($\Leftarrow$) If $a\in\Hit(C)$, choose $x\in C$ with $x\forces a$. Since $x\in\Sel(U)$ we have $a\in N(x)\subseteq U$.
\end{proof}

\begin{remark}[Meaning]
$\Hit(C)$ records which generators are \emph{witnessed} in $C$.
$\Sel(U)$ picks points whose \emph{entire neighborhood behavior} is constrained by $U$.
This is the ``closed-without-complements'' move: $U$ acts as a \emph{code} for a closed-like object by constraining neighborhoods, not by negation.
\end{remark}

% =========================================================
\section{Positivity and closed sets without complements}

\subsection{The positivity operator $J=\Hit\circ\Sel$}

\begin{definition}[Positivity interior on codes]
Define
\[
J:\Pow(S)\to\Pow(S),\qquad J(U):=\Hit(\Sel(U)).
\]
Equivalently,
\[
a\in J(U)\quad\Longleftrightarrow\quad \exists x\in X\ \bigl(x\forces a\ \wedge\ N(x)\subseteq U\bigr).
\]
\end{definition}

\begin{definition}[Positivity relation]
Define $\ltimesrel\subseteq S\times \Pow(S)$ by
\[
a\ltimesrel U \quad:\Longleftrightarrow\quad a\in J(U).
\]
\end{definition}

\begin{remark}[One-line intuition]
$a\ltimesrel U$ means: there exists a point $x$ in basic open $a$ such that all basic neighborhoods of $x$ lie in $U$.
So $U$ behaves like a ``closed neighborhood specification'' and $a$ \emph{hits} it.
\end{remark}

\begin{example}[Medical: Positivity as witnessed symptom clusters]\label{ex:medical-positivity}
The statement $\mathsf{fever} \ltimesrel \{\mathsf{fever}, \mathsf{cough}, \mathsf{fatigue}\}$ means:

\begin{quote}
There exists a patient who has fever and whose \emph{complete} symptom profile is $\subseteq \{\mathsf{fever}, \mathsf{cough}, \mathsf{fatigue}\}$.
\end{quote}

This patient is a \textbf{witness} that fever is consistent with the symptom cluster.

Contrast: $\mathsf{fever} \ltimesrel \{\mathsf{hypothermia}\}$ is \emph{false} because no patient
has both fever and only hypothermia --- the combination is contradictory.

Positivity detects \textbf{realizability}: which symptom combinations actually occur in patients.
\end{example}

\begin{proposition}[A point forces positivity]\label{prop:point-implies-positive}
Assume a point $x\in X$ and a generator $a\in S$ with $x\forces a$.
Then $a$ is positive in the unary sense, i.e.\ $a\ltimesrel S$.
More generally, if $U\subseteq S$ satisfies $N(x)\subseteq U$, then $a\ltimesrel U$.
\end{proposition}

\begin{proof}
If $N(x)\subseteq U$ then $x\in\Sel(U)$ by definition, hence $a\in\Hit(\Sel(U))=J(U)$,
i.e.\ $a\ltimesrel U$.  Taking $U=S$ gives $a\ltimesrel S$.
\end{proof}

\subsection{Elementary properties}

\begin{proposition}[Contractive and monotone]
For all $U,V\subseteq S$:
\begin{enumerate}[label=(\alph*)]
\item $J(U)\subseteq U$.
\item If $U\subseteq V$ then $J(U)\subseteq J(V)$.
\end{enumerate}
\end{proposition}

\begin{proof}
(a) If $a\in J(U)$, choose $x$ with $x\forces a$ and $N(x)\subseteq U$. Then $a\in N(x)\subseteq U$.

(b) If $U\subseteq V$ then $\Sel(U)\subseteq \Sel(V)$, hence $\Hit(\Sel(U))\subseteq \Hit(\Sel(V))$.
\end{proof}

\begin{proposition}[Idempotence]\label{prop:idempotence-laws}
$J= \Hit\circ\Sel$ is idempotent: $J(J(U))=J(U)$.
\end{proposition}

\begin{proof}
Since $J\le \mathrm{id}$, we have $J(J(U))\subseteq J(U)$.
For the reverse, $\mathrm{id}\le \Sel\circ\Hit$ implies $\Sel(U)\subseteq \Sel(J(U))$; applying $\Hit$ yields $J(U)\subseteq J(J(U))$.
\end{proof}

\begin{definition}[Formal closed (via positivity)]
Call $U\subseteq S$ \emph{(positively) closed} if $J(U)=U$.
\end{definition}

\begin{example}[Medical: Complete symptom profiles]\label{ex:medical-closed}
A set $U \subseteq S$ is positively closed iff: whenever there exists a patient with symptom $a$
whose profile is contained in $U$, then $a \in U$.

Equivalently: $U$ is a \textbf{complete symptom cluster} --- you cannot find a patient ``inside $U$''
who exhibits a symptom outside $U$.
\end{example}

\subsection{The compatibility axiom}

The key interaction between covers and positivity is captured by:

\begin{proposition}[Compatibility]\label{prop:compatibility}
For all $a \in S$ and $U, V \subseteq S$:
\[
(a \covers U) \land (a \ltimesrel V) \implies \exists u \in U.\; u \ltimesrel V.
\]
\end{proposition}

\begin{proof}
Suppose $a \covers U$ and $a \ltimesrel V$.
By positivity, there exists $x$ with $x \forces a$ and $N(x) \subseteq V$.
By the cover, since $x \forces a$, there exists $u \in U$ with $x \forces u$.
Then $u \in N(x) \subseteq V$ and $x$ witnesses $u \ltimesrel V$.
\end{proof}

\begin{example}[Medical: Refinement preserves witnesses]
Suppose $\mathsf{pneumonia} \covers \{\mathsf{fever}, \mathsf{cough}, \mathsf{infiltrates}\}$
and there exists a pneumonia patient confined to profile $V$.

Then there exists a fever patient (or cough patient, or infiltrates patient) also confined to $V$.

The diagnostic refinement from ``pneumonia'' to its symptoms preserves the witness.
\end{example}

\begin{remark}[Why compatibility matters]
Compatibility is the \textbf{soundness condition} for refinement:
if you refine a hypothesis (via covers), you don't lose your witnesses (positivity).
This is what makes the framework coherent for inference.
\end{remark}

% =========================================================
\section{Information content and refinement games}\label{sec:info-games}

The formal relations $\covers$ and $\ltimesrel$ admit an interpretation closer to computation than to classical topology.
Think of a generator $a\in S$ as a \emph{partial information state}.  A cover $a\covers U$
says that to \emph{settle} the information in $a$ one must proceed to a feasible refinement and end up in \emph{one of} the subcases in $U$.

Positivity adds one crucial twist: existence is not classical negation of emptiness, but a \emph{witnessed} possibility of continuing the refinement process.

\subsection{The refinement game $G(a,U)$}

Fix $a\in S$ and a set $U\subseteq S$ of ``acceptable outcomes''.
Define an infinite two-player game $G(a,U)$:

\begin{itemize}
  \item The current position is a generator $b\in S$ (initially $b=a$).
  \item \emph{Refiner} (``Devil'') chooses a cover $b\covers V$.
  \item \emph{Witness} (``God'') responds by choosing some $v\in V$.
  \item The next position is $b:=v$.
  \item Witness wins if she can play forever while keeping every chosen position inside $U$.
\end{itemize}

\begin{proposition}[Positivity as a winning condition]\label{prop:positivity-game}
If $a\ltimesrel U$, then Witness has a winning strategy in $G(a,U)$.
\end{proposition}

\begin{proof}
By $a\in J(U)$ there exists $x\in X$ with $x\forces a$ and $N(x)\subseteq U$.
Whenever Refiner plays $b\covers V$ with $x\forces b$, there exists $v\in V$ with $x\forces v$.
Witness chooses such $v$. Since $v\in N(x)\subseteq U$, the position stays in $U$.
\end{proof}

\paragraph{Medical diagnosis as a refinement game.}
The game $G(\mathsf{fever}, U)$ where $U$ is a diagnostic profile:
\begin{itemize}
  \item Refiner (skeptic): ``You say fever? What kind? Bacterial? Viral?''
  \item Witness (diagnostician): ``Bacterial fever'' (chooses refinement in $U$)
  \item Refiner: ``What symptoms? Give me a cover.''
  \item Witness: ``Elevated WBC'' (still in $U$)
\end{itemize}
Witness wins iff there's a real patient (the witness $x$) to guide all choices.

\paragraph{Legal reasoning as a refinement game.}
Let $X$ = legal cases, $S$ = claims/evidence/precedents.
The game: opposing counsel (Refiner) demands finer distinctions; the advocate (Witness) must find precedents.
Positivity = having a real case to cite at every challenge.

% =========================================================
\section{Two adjunctions and four operators}\label{sec:two-adjunctions}

The same forcing table gives two adjunctions, hence four operators:

\[
\Pow(S)\ \underset{\Intt}{\overset{\Ext}{\rightleftarrows}}\ \Pow(X)
\qquad(\Ext\dashv \Intt),
\qquad\qquad
\Pow(X)\ \underset{\Sel}{\overset{\Hit}{\rightleftarrows}}\ \Pow(S)
\qquad(\Hit\dashv \Sel).
\]

\begin{center}
\begin{tabular}{lll}
\toprule
\textbf{Operator} & \textbf{Type} & \textbf{Meaning} \\
\midrule
$\Ext$ & $\Pow(S) \to \Pow(X)$ & Points satisfying some observable in $U$ \\
$\Intt$ & $\Pow(X) \to \Pow(S)$ & Observables universally valid on $A$ \\
$\Hit$ & $\Pow(X) \to \Pow(S)$ & Observables witnessed in $C$ \\
$\Sel$ & $\Pow(S) \to \Pow(X)$ & Points confined to profile $U$ \\
\midrule
$j = \Intt \circ \Ext$ & $\Pow(S) \to \Pow(S)$ & Closure (formal opens) \\
$c = \Ext \circ \Intt$ & $\Pow(X) \to \Pow(X)$ & Interior (extensional opens) \\
$J = \Hit \circ \Sel$ & $\Pow(S) \to \Pow(S)$ & Interior (formal closeds) \\
$\cl = \Sel \circ \Hit$ & $\Pow(X) \to \Pow(X)$ & Closure (on points) \\
\bottomrule
\end{tabular}
\end{center}

\begin{remark}[Two sides of the same coin]
Both adjunctions arise from the same forcing relation, but govern different semantics:
\begin{itemize}
\item $\Ext\dashv\Intt$ = universal refinement / locale / covers
\item $\Hit\dashv\Sel$ = existential probes / positivity / closed-without-complements
\end{itemize}
\end{remark}

% =========================================================
\section{Reconstruction of forcing from residuals}

A central observation (cf.\ Ciraulo--Sambin~\cite{CirauloSambin2012}) is that the forcing table is recoverable from either left adjoint,
since left adjoints preserve joins and are determined by their action on singletons.

\begin{theorem}[Singleton reconstruction]
If $F=\Ext$ arises from $\forces$, then $\forces$ is uniquely determined by:
\[
x\forces a \Longleftrightarrow x\in \Ext(\{a\}).
\]
Similarly for $\Hit$:
\[
x\forces a \Longleftrightarrow a\in \Hit(\{x\}).
\]
\end{theorem}

\begin{proof}
Since $\Ext$ preserves unions, $\Ext(U)=\bigcup_{a\in U}\Ext(\{a\})$.
Thus $x\in \Ext(U)$ iff $\exists a\in U$ with $x\in \Ext(\{a\})$.
\end{proof}

\begin{remark}[What is \emph{not} reconstructible]
The composites $j=\Intt\circ\Ext$ or $J=\Hit\circ\Sel$ generally lose witness-level information.
The residual \emph{pair} determines forcing; a single composite does not.
\end{remark}

% =========================================================
\section{Costs and resource-bounded semantics}

The Boolean case uses $\{0,1\}$. We now generalize by replacing Boolean truth with values in an ordered cost structure,
mirroring the enriched positivity relations studied by Ciraulo--Vickers~\cite{CirauloVickers2016}.

\subsection{Ordered costs: a minimal baseline}\label{subsec:ordered-costs}

\begin{definition}[Cost set]\label{def:cost-set}
A \emph{cost set} is a preorder $(B,\preceq)$ whose elements are budgets.
We read $b\preceq b'$ as ``budget $b$ is usable whenever $b'$ is'' (more budget = more permissive).
\end{definition}

\begin{definition}[Cost-valued forcing]\label{def:cost-valued-forcing}
A \emph{cost-valued forcing matrix} is a map $R:X\times S\to B$.
For each budget $b\in B$ we define:
\[
x \forces_b a \quad:\!\!\iff\quad R(x,a)\preceq b.
\]
\end{definition}

\begin{example}[Medical: Verification costs]\label{ex:medical-costs}
Different tests have different costs:

\begin{center}
\begin{tabular}{lrl}
\toprule
\textbf{Test/Observation} & \textbf{Cost} & \textbf{Notes} \\
\midrule
Fever & \$0 & Thermometer \\
Cough & \$0 & Observation \\
Blood pressure & \$5 & Standard equipment \\
Blood test (WBC) & \$50 & Lab work \\
X-ray & \$200 & Imaging \\
CT scan & \$500 & Advanced imaging \\
PCR test & \$100 & Molecular \\
\bottomrule
\end{tabular}
\end{center}

The cost-valued forcing $R(x, a) = c$ means: ``verifying that patient $x$ has condition $a$ costs $c$.''

At budget $b = \$50$, we can verify fever, cough, blood pressure, and WBC --- but not X-ray or CT scan.

The resource-aware diagnostic question: \emph{``What is the cheapest test that discriminates between remaining hypotheses?''}
\end{example}

\begin{remark}[Filtration viewpoint]
For $b\preceq b'$ we have $x\forces_b a \Rightarrow x\forces_{b'} a$.
All constructions from earlier sections apply at each budget $b$, producing a monotone filtration.
\end{remark}

\subsection{Budget-indexed covers and positivity}

\begin{definition}[Budget-indexed relations]
For each budget $b \in B$:
\[
a \covers_b U \quad:\Longleftrightarrow\quad \forall x,\ (x \forces_b a \Rightarrow \exists u \in U,\ x \forces_b u)
\]
\[
a \ltimesrel_b U \quad:\Longleftrightarrow\quad \exists x,\ (x \forces_b a \land N_b(x) \subseteq U)
\]
where $N_b(x) = \{a \in S : x \forces_b a\}$.
\end{definition}

\begin{example}[Medical: Budget-dependent diagnosis]
At budget \$50 (fever, cough, blood test available):
\begin{itemize}
  \item $\mathsf{bacterial\_infection} \covers_b \{\mathsf{fever}, \mathsf{elevated\_WBC}\}$ --- can verify
  \item $\mathsf{pneumonia} \covers_b \{\mathsf{fever}, \mathsf{cough}\}$ --- partial (can't see infiltrates)
\end{itemize}

At budget \$250 (adds X-ray):
\begin{itemize}
  \item $\mathsf{pneumonia} \covers_{b'} \{\mathsf{fever}, \mathsf{cough}, \mathsf{infiltrates}\}$ --- full verification
\end{itemize}

Higher budget = finer discrimination = more refined diagnosis.
\end{example}

\subsection{Optional: cost domains}

\begin{definition}[Cost domain]\label{def:cost-domain}
A \emph{cost domain} is $(V,\le,\otimes,I,\multimap)$ where:
\begin{enumerate}[label=(\alph*)]
\item $(V,\le)$ is a poset
\item $(V,\otimes,I)$ is a commutative monoid, monotone in both arguments
\item Residuation: $u\otimes v \le w \Longleftrightarrow v \le (u\multimap w)$
\item Specified joins/meets exist for the families used below
\end{enumerate}
\end{definition}

\begin{remark}[Examples]
\begin{enumerate}[label=(\roman*)]
\item Boolean: $V=\{0,1\}$, $\otimes=\wedge$, $u\multimap w=(\neg u)\vee w$
\item Lawvere metric: $V=[0,\infty]$ reversed, $\otimes=+$, $u\multimap w=\max\{0,w-u\}$
\item Probability: $V = [0,1]$, $\otimes = \cdot$, $I = 1$
\end{enumerate}
\end{remark}

\subsection{Enriched adjunctions}

With a cost domain $V$ and $V$-valued forcing $R:X\times S\to V$:

\begin{definition}[Enriched operators]
\[
(\Ext_V\phi)(x)=\bigvee_{a\in S}\ \phi(a)\otimes R(x,a),
\qquad
(\Int_V\psi)(a)=\bigwedge_{x\in X}\ \bigl(R(x,a)\multimap \psi(x)\bigr)
\]
\[
(\Hit_V\psi)(a)=\bigvee_{x\in X}\ \psi(x)\otimes R(x,a),
\qquad
(\Sel_V\phi)(x)=\bigwedge_{a\in S}\ \bigl(R(x,a)\multimap \phi(a)\bigr)
\]
\end{definition}

\begin{proposition}[Enriched adjunctions]
$\Ext_V\dashv \Int_V$ and $\Hit_V\dashv \Sel_V$ (when joins/meets exist).
\end{proposition}

\begin{definition}[Enriched positivity]
$J_V:=\Hit_V\circ \Sel_V$ and $\|a\ltimesrel \phi\|_V := (J_V\phi)(a)$.
\end{definition}

\begin{remark}
When $V=\{0,1\}$, this recovers the Boolean case exactly.
\end{remark}

% =========================================================
\section{Conceptual summary: the three-layer separation}

\begin{center}
\fbox{\begin{minipage}{0.9\linewidth}
\textbf{Forcing-matrix viewpoint (three-layer separation).}
\[
\begin{array}{c}
\text{(i) forcing/adjoints: } R:X\times S\to V\\[2pt]
\Downarrow\\[2pt]
\text{(ii) positivity \& closedness without complements: } J, \ltimesrel\\[2pt]
\Downarrow\\[2pt]
\text{(iii) resource-aware logic: behavior of }\wedge\text{ vs.\ }\otimes\text{, structural rules}
\end{array}
\]
Layer (iii) is where resource-bounded positive topology deviates from the classical picture.
\end{minipage}}
\end{center}

\subsection{Budget-indexed truth}

The syntax remains $(\wedge,\vee,\to,\forall,\exists)$.
What changes is the judgement: truth is indexed by budget.
\[
a \Vdash_b \varphi \quad\text{means: ``from state $a$, $\varphi$ is verified within budget $b$''}
\]

\paragraph{Why $\wedge$ becomes subtle.}
To establish $\varphi\wedge\psi$ you must allocate resources across subgoals.
We use an abstract splitting relation $\Split(b;\,b_1,b_2)$:
\[
a\Vdash_b(\varphi\wedge\psi)\quad\Longleftrightarrow\quad
\exists b_1,b_2\ \bigl(\Split(b;\,b_1,b_2)\ \&\ a\Vdash_{b_1}\varphi\ \&\ a\Vdash_{b_2}\psi\bigr)
\]

\paragraph{Existence is witnessed.}
To assert $\exists x.\,\varphi(x)$ one must provide a witness fitting the budget.
Positivity is the mechanism turning ``nonempty'' into ``witnessably inhabited''.

% =========================================================
\appendix
\section{Application to AI engineering}

This appendix explains why forcing/positivity/cost provides concrete engineering value for AI systems.

\subsection{The AI analogue of forcing}

In AI engineering, one evaluates systems against requirements:
\begin{itemize}
\item $X$ = executions / traces / agent rollouts
\item $S$ = constraints (``tool-call succeeded'', ``no disallowed content'', ``latency $<500$ms'', etc.)
\item $x\forces a$ = ``execution $x$ satisfies constraint $a$''
\end{itemize}

The forcing matrix $R:X\times S\to V$ is an evaluation matrix.

\subsection{Two adjunctions = two engineering operations}

\paragraph{(1) $\Ext\dashv\Intt$: specification refinement.}
\begin{itemize}
\item $\Ext(U)$: which traces satisfy at least one check in $U$?
\item $\Intt(A)$: which checks are universally valid over traces $A$?
\item $a\covers U$: whenever $a$ holds, some $b\in U$ holds
\end{itemize}
This is specification management: what the system guarantees.

\paragraph{(2) $\Hit\dashv\Sel$: witness finding.}
\begin{itemize}
\item $\Hit(C)$: which checks are witnessed in trace-set $C$?
\item $\Sel(U)$: which traces live entirely inside constraint envelope $U$?
\item $a\ltimesrel U$: exists a trace satisfying $a$ with profile $\subseteq U$
\end{itemize}
This is counterexample search and safe operating region identification.

\subsection{Positivity = safety envelopes}

Engineers want ``safe = complement of bad''. This fails because:
\begin{enumerate}[label=(\alph*)]
\item the complement is not observable (unknown unknowns)
\item violations are sparse and adversarial
\item safety is better defined by positive evidence of robustness
\end{enumerate}

Positivity gives a constructive alternative:
\begin{itemize}
\item closed code $U$ with $J(U)=U$ is a safety envelope
\item $a\ltimesrel U$ provides witness-based membership
\end{itemize}

Safety is a robust region, not a complement.

\subsection{LLM + PosTop: grounded generation}\label{subsec:llm-postop}

Current LLM agents fail in three ways:
\begin{enumerate}
  \item \textbf{Hallucination}: claims without witnesses
  \item \textbf{Contradiction}: inconsistent claims
  \item \textbf{Inefficiency}: expensive verification when cheap suffices
\end{enumerate}

PosTop addresses all three:
\begin{itemize}
  \item \textbf{Cover check}: Does conclusion follow from observations?
  \item \textbf{Positivity check}: Is there a witness?
  \item \textbf{Resource optimization}: What's the cheapest discriminating test?
\end{itemize}

\paragraph{Architecture.}
\begin{enumerate}
  \item LLM generates candidate hypotheses $H$
  \item PosTop filters: keep $h \in H$ iff $h \covers \text{observed}$ and $h \ltimesrel \text{observed}$
  \item PosTop suggests: cheapest observation to discriminate remaining $H$
  \item Repeat until $|H| = 1$ or budget exhausted
\end{enumerate}

This turns the LLM from oracle to \emph{grounded reasoner}: every claim has a cover-chain (explanation) and witness (grounding).

\paragraph{Comparison with knowledge graphs.}

\begin{center}
\begin{tabular}{lcc}
\toprule
\textbf{Feature} & \textbf{Knowledge Graph} & \textbf{PosTop} \\
\midrule
Entities & \checkmark & \checkmark \\
Relations & \checkmark & \checkmark \\
Facts & \checkmark & \checkmark \\
\textbf{Entailment} & $\times$ & \checkmark \\
\textbf{Witnesses} & $\times$ & \checkmark \\
\textbf{Costs} & $\times$ & \checkmark \\
\textbf{Refinement soundness} & $\times$ & \checkmark \\
\bottomrule
\end{tabular}
\end{center}

\begin{quote}
\textbf{Knowledge Graphs} tell you what exists.\\
\textbf{PosTop} tells you what's consistent, what implies what, and how to verify efficiently.
\end{quote}

Knowledge graphs are dictionaries. PosTop is grammar.

% =========================================================
\appendix
\section{Software and Reproducibility}\label{app:software}

The LaTeX file you are reading lives in the same repository as the reference implementation of the forcing matrix, adjunctions, and notebooks.
The codebase is organized as follows:
\begin{itemize}
  \item \texttt{src/postop/core.py} implements the data structures discussed in Sections~\ref{sec:adj2}--\ref{sec:info-games}, with the forcing relation encoded as a Python dictionary of point $\to$ observable sets, plus the cost-aware variants in \texttt{src/postop/costs.py}.
  \item \texttt{src/postop/operators.py} and \texttt{src/postop/llm.py} provide tracing/explainability helpers that surface the cover witnesses/counterexamples required by the diagnostic examples.
  \item \texttt{examples/} mirrors the running medical and guardrail scenarios; \texttt{tests/} exercises the adjunction identities via \texttt{pytest}.
  \item \texttt{docs/REPRODUCIBILITY.md} gives a copy/paste runbook matching the commands below.
\end{itemize}
The latest sources live at \url{https://github.com/mirco-mannucci/postop}; we snapshot paper+code together via the git tag \texttt{v0.1.0} (and subsequent releases) so the commands below always map to a reproducible commit/DOI pair.

\subsection{Environment setup}

We keep the environment lightweight (pure Python + pytest). From the repository root:
\begin{verbatim}
python3 -m venv .venv
source .venv/bin/activate
pip install -e .[dev]
\end{verbatim}

\subsection{Running the examples}

The forcing table is serialized as ordinary Python dictionaries and can be inspected directly. To reproduce the medical and guardrail walkthroughs from Sections~\ref{sec:adj2} and~\ref{sec:info-games}:
\begin{verbatim}
python examples/medical.py
python examples/llm_guardrail.py
\end{verbatim}
Each script prints the Ext/Int/Hit/Sel results, the compatibility witnesses, and the guardrail counterexamples that correspond to the formal statements in the main text.

\subsection{Notebooks and tests}

The ActionList roadmap adds four short notebooks under \texttt{notebooks/}; they can be executed headlessly via:
\begin{verbatim}
jupyter nbconvert --to notebook --execute notebooks/01_basics_ext_int_hit_sel_j.ipynb
\end{verbatim}
Swap in the remaining notebook filenames (medical diagnosis, costs, guardrails) to reproduce the figures in the extended version of the paper.

Finally, run the test suite to check the adjunction/proof obligations:
\begin{verbatim}
pytest -q
\end{verbatim}
These tests import the installed package so they are venv-friendly and catch regressions both in the mathematical invariants (idempotence/contractivity) and in the explainability helpers referenced by the notebooks.

% =========================================================
\begin{thebibliography}{99}

\bibitem{Sambin-OUP-PositiveTopology}
G.~Sambin, \emph{Positive Topology: A New Practice in Constructive Mathematics}.
Oxford University Press, 2025.

\bibitem{CirauloSambin2012}
F.~Ciraulo and G.~Sambin.
A constructive {G}alois connection between closure and interior.
\emph{J. Symbolic Logic} 77(4):1308--1324, 2012.

\bibitem{Valentini2012}
S.~Valentini.
Relative formal topology: the binary positivity predicate comes first.
\emph{Math. Structures Comput. Sci.} 22(1):69--102, 2012.

\bibitem{CirauloVickers2016}
F.~Ciraulo and S.~Vickers.
Positivity relations on a locale.
\emph{Ann. Pure Appl. Logic} 167(9):806--819, 2016.

\bibitem{MacLaneMoerdijk1992}
S.~Mac~Lane and I.~Moerdijk.
\emph{Sheaves in Geometry and Logic}.
Springer, 1992.

\bibitem{CoquandSambinSmithValentini2003}
T.~Coquand, G.~Sambin, J.~Smith, and S.~Valentini.
Inductively generated formal topologies.
\emph{Ann. Pure Appl. Logic} 124(1--3):71--106, 2003.

\end{thebibliography}

\end{document}
